\Article{5}
Der Kürze halber setzen wir
\[
\begin{aligned}
(aa+bb+cc)\mathrm{d}t^2+2(aa'+bb'+cc')\mathrm{d}t\mathord{.}\mathrm{d}u&\\
+(a'a'+b'b'+c'c')\mathrm{d}u^2&=\omega
\end{aligned}
\]
\noindent und bemerken, dass die Differentialgleichung \(\omega=0\) \textit{zwei} Integrationen zulässt. Zerlegt man nämlich das Trinom \(\omega\) in zwei Faktoren, die in
\SourceLeaf{201}
\(\mathrm{d}t\) und \(\mathrm{d}u\) linear sind, so muss entweder der eine oder der andere Faktor \(=0\) werden, was zwei verschiedene Integrationen ergibt. Die eine Integration wird der Gleichung
\[
\begin{aligned}
0={}&(aa+bb+cc)\mathrm{d}t\\
&+\bigl\{aa'+bb'+cc'+i\BareRoot{\bigl((aa+bb+cc)(a'a'+b'b'+c'c')}\\
&\qquad-(aa'+bb'+cc')^2\bigr)\bigr\}\mathrm{d}u
\end{aligned}
\]
\noindent entsprechen (wobei \(i\) der Kürze halber für \(\BareRoot{-1}\) geschrieben ist, da man sich leicht davon überzeugt, dass der irrationale Teil des Ausdrucks imaginär werden muss); die andere entspricht einer ganz ähnlichen Gleichung, in der nur \(i\) durch \(-i\) ersetzt wird. Lautet also das Integral der ersten Gleichung
\[
p+iq=\text{Const.}
\]
\noindent wobei \(p\) und \(q\) reelle Funktionen von \(t\) und \(u\) bezeichnen, so lautet das andere Integral
\[
p-iq=\text{Const.}
\]
\noindent und aus der Natur der Sache folgt, dass
\[
(\mathrm{d}p+i\,\mathrm{d}q)\mathord{.}(\mathrm{d}p-i\,\mathrm{d}q)\quad\text{oder}\quad\mathrm{d}p^2+\mathrm{d}q^2
\]
\noindent ein Faktor von \(\omega\) sein muss, oder dass
\[
\omega=n(\mathrm{d}p^2+\mathrm{d}q^2)
\]
\noindent gelten muss, wobei \(n\) eine endliche Funktion von \(t\) und \(u\) ist.

Wir wollen nun das Trinom, in das
\[
\mathrm{d}X^2+\mathrm{d}Y^2+\mathrm{d}Z^2
\]
\noindent übergeht, wenn für \(\mathrm{d}X,\mathrm{d}Y,\mathrm{d}Z\) die durch \(T,U,\mathrm{d}T,\mathrm{d}U\) ausgedrückten Werte eingesetzt werden, mit \(\Omega\) bezeichnen und annehmen, dass wie zuvor die beiden Integrale der Gleichung \(\Omega=0\) lauten:
\[
\begin{aligned}
P+iQ&=\text{Const.}\\
P-iQ&=\text{Const.}
\end{aligned}
\]
\noindent und
\[
\Omega=N(\mathrm{d}P^2+\mathrm{d}Q^2)
\]
\noindent wobei \(P,Q,N\) reelle Funktionen von \(T\) und \(U\) bezeichnen.
\par\indent\SourceLeaf{202}%
Diese Integrationen lassen sich (wenn man von den allgemeinen Schwierigkeiten der Integration absieht) offenbar vor der Lösung unserer Hauptaufgabe ausführen.

Werden nun für \(T,U\) solche Funktionen von \(t,u\) eingesetzt, die die Bedingung unserer Hauptaufgabe erfüllen, so geht \(\Omega\) in \(mm\omega\) über, und es ergibt sich
\[
\frac{(\mathrm{d}P+i\,\mathrm{d}Q)\mathord{.}(\mathrm{d}P-i\,\mathrm{d}Q)}{(\mathrm{d}p+i\,\mathrm{d}q)\mathord{.}(\mathrm{d}p-i\,\mathrm{d}q)}=\frac{mmn}{N}
\]
\noindent Man sieht aber leicht, dass der Zähler auf der ersten Seite dieser Gleichung nur dann durch den Nenner teilbar sein kann, wenn
\[
\begin{aligned}
&\text{entweder }\mathrm{d}P+i\,\mathrm{d}Q\text{ durch }\mathrm{d}p+i\,\mathrm{d}q,\text{ und }\mathrm{d}P-i\,\mathrm{d}Q\text{ durch }\mathrm{d}p-i\,\mathrm{d}q,\\
&\text{oder }\mathrm{d}P+i\,\mathrm{d}Q\text{ durch }\mathrm{d}p-i\,\mathrm{d}q,\text{ und }\mathrm{d}P-i\,\mathrm{d}Q\text{ durch }\mathrm{d}p+i\,\mathrm{d}q
\end{aligned}
\]
\noindent die angegebenen Teilbarkeiten gelten. Im ersten Fall verschwindet demnach \(\mathrm{d}P+i\,\mathrm{d}Q\), wenn \(\mathrm{d}p+i\,\mathrm{d}q=0\), oder \(P+iQ\) wird konstant, wenn \(p+iq\) als konstant angenommen wird; das heißt, \(P+iQ\) ist nur eine Funktion von \(p+iq\) und ebenso \(P-iQ\) eine Funktion von \(p-iq\). Im anderen Fall ist \(P+iQ\) eine Funktion von \(p-iq\) und \(P-iQ\) eine Funktion von \(p+iq\). Es ist leicht einzusehen, dass diese Folgerungen auch umgekehrt gelten, nämlich dass, wenn für \(P+iQ,P-iQ\) Funktionen von \(p+iq,p-iq\) (entweder jeweils in der angegebenen oder in umgekehrter Reihenfolge) angenommen werden. die endliche Teilbarkeit von \(\Omega\) durch \(\omega\) und damit die oben als notwendig erkannte Proportionalität vorliegt.

Man überzeugt sich außerdem leicht davon, dass, wenn zum Beispiel
\[
\begin{aligned}
P+iQ&=f(p+iq)\\
P-iQ&=f'(p-iq)
\end{aligned}
\]
\noindent gesetzt werden, die Beschaffenheit der Funktion \(f'\) schon durch die von \(f\) bestimmt ist. Wenn nämlich unter den konstanten Größen, die diese enthalten kann, nur reelle vorkommen, muss die andere Funktion \(f'\) mit \(f\) ganz identisch sein, damit reellen Werten von \(p,q\) stets reelle Werte von \(P,Q\) entsprechen. Andernfalls unterscheidet sich \(f'\) von \(f\) nur darin, dass in den imaginären Elementen von \(f\) anstelle von \(i\) überall das entgegengesetzte \(-i\) gesetzt werden muss.

Man erhält sodann
\SourceLeaf{203}
\[
\begin{aligned}
P&=\tfrac{1}{2}f(p+iq)+\tfrac{1}{2}f'(p-iq)\\
iQ&=\tfrac{1}{2}f(p+iq)-\tfrac{1}{2}f'(p-iq)
\end{aligned}
\]
\noindent oder, was dasselbe ist: Wenn die Funktion \(f\) ganz beliebig gewählt wird (nach Wunsch auch mit konstanten imaginären Elementen), setzt man \(P\) gleich dem reellen Teil und \(iQ\) (bei der zweiten Lösung \(-iQ\)) gleich dem imaginären Teil von \(f(p+iq)\). Daraus werden dann durch Elimination \(T\) und \(U\) als Funktionen von \(t\) und \(u\) dargestellt. Damit ist die gestellte Aufgabe ganz allgemein und vollständig gelöst.

